\chapter{Basic operations}

\section{Addition, subtraction, multiplication and division}

\noindent Here we present the most basic mathematical operations on the real numbers. We will only mention shortly the main properties as a reminder.

\subsection{Basic operations on real numbers}

\subsubsection{Addition}

\noindent Take any two real numbers $a$ and $b$. It is true no matter what $a$ and $b$ are that
$$
a + b = b + a
$$
Mathematicians call this property commutativity and say that the sum of two real numbers is commutative. But that is just a fancy word. And we should not focus on names. Basically, the order of adding two numbers does not affect the result.

\subsubsection{Subtraction}

\noindent We remain with our numbers $a$ and $b$. By $-b$ we mean the number you get when changing $b$'s sign. Notice that $-b$ isn't necessarily negative. If $b$ itself is negative, changing its sign will make it positive and so $-b$ will be positive in such a case.\\

\noindent The difference of $a$ and $b$ is then defined as
$$
a - b = a + (-b)
$$
In words, take the numbers $a$ and $b$, change the sign of $b$ and add $-b$ to $a$.\\

\noindent Because the sum is commutative it follows that
$$
a - b = a + (-b) = -b + a
$$
But unlike addition, the order of the numbers in a subtraction cannot be reversed without changing the result:
$$
a - b \neq b - a
$$
Changing the order of the subtraction amounts to changing the sign of the result:
$$
a - b = -(b-a)
$$

\subsubsection{Multiplication}

\noindent Writing a sum can be tedious:
\begin{equation*}
5 + 5 + 5 + 5 + 5 + 5 + 5 = 35
\end{equation*}
A multiplication is in the simplest case a shorthand notation for this. The number 5 is added 7 times here, which is written as: $ 5 \times 7 $.\\

\noindent As was the case for the addition, the multiplication is commutative:
\begin{equation*}
4 \times 8 = 8 \times 4
\end{equation*}
This equality says that $8+8+8+8 = 4+4+4+4+4+4+4+4$, which is true.\\

\noindent Multiplying 2 positive numbers or 2 negative numbers yields a positive result. When one of the numbers is positive and the other is negative the result is negative. And multiplying any number by 0 always gives 0.
\begin{eqnarray*}
-5 \times 9 & = & -45 \\
17 \times 0 & = & 0 \\
-6 \times (-8) & = & 48
\end{eqnarray*}
\noindent As we will soon get to equations and variables in our discussion of mathematics, it may get confusing if we use the symbol $\times$ for multiplication and $x$ for a variable. In order to avoid confusion we will often write multiplications with a "." symbol: $3 \times 4 = 3 . 4$. The symbols "$.$" and "$\times$" both represent the same action of multiplying the numbers on either side of them with each other, and they will be used interchangeably throughout the rest of the text. Additionally, when representing numbers by letters, we will often omit any sign of multiplication all together: $a \times b = a . b = ab$. They all mean $a$ times $b$. We have actually already done this earlier in the text when we discussed properties of rational numbers in chapter \ref{Set_Theory}. We wrote there $1\: 000 \: x$ which means one thousand times $x$, though we omitted and sign of multiplication.

\subsubsection{Division}

\noindent The division is the inverse function of the multiplication. This means that if $a \times b = c$ then $a = \frac{c}{b}$ or equally $b = \frac{c}{a}$.
$$
\frac{63}{7} = 9 \text{ is equivalent with } 7 \times 9 = 63
$$
\noindent Division of $a$ and $b$ is written as $\frac{a}{b}$ or as $a \div b$.\\

\noindent \textbf{Point of attention:} The denominator in a division is never allowed to be zero. We will see when discussing the concept of limits how cases that look like a division by zero can be handled. But as such, dividing by zero is a very big no-no in mathematics.

\subsection{Basic operations of fractions}

\noindent A fraction is nothing more than a division that still needs to be done. It is of the form $\frac{a}{b}$ where $a$ is called the numerator and $b$ the denominator. The denominator has to be different from 0.

\subsubsection{Simplifying fractions}

\noindent Two fractions are equal when the result of the division of numerator by denominator is equal:
\begin{equation*}
\frac{1}{2} = 0.5 = \frac{8}{16}
\end{equation*}
A fraction remains unchanged when multiplying or dividing both numerator and denominator by the same number as this operation does not change the final result of their division:
\begin{equation*}
\frac{4}{15} = \frac{4 \times 3}{15 \times 3} = \frac{12}{45}
\end{equation*}
As you can easily check these 2 fractions indeed yield the same result when dividing numerator by denominator. Simplifying a fractions means getting rid of all the common factors in numerator and denominator. The fraction $\frac{144}{20}$ can be simplified to $\frac{36}{5}$. To do this we write it as
\begin{eqnarray*}
\frac{144}{20} = \frac{4 \times 36}{4 \times 5}
\end{eqnarray*}
\noindent By canceling out the $4$, which is a common factor in the numerator and denominator, we are indeed left with $\frac{36}{5}$ which does not have any remaining common factors.\\

\subsubsection{Addition and subtraction of fractions}

\noindent When adding (subtracting) fractions we need to make sure they have the same denominator. When the denominators are the same, the sum (difference) of the fractions is the sum (difference) of the numerators divided by their common denominator:
\begin{eqnarray*}
\frac{5}{2} + \frac{17}{2} = \frac{22}{2} = \frac{11}{1} = 11\\
-\frac{33}{16} + \frac{1}{16} = -\frac{32}{16} = -\frac{2}{1} = -2
\end{eqnarray*} 
When the denominators are unequal, we need to multiply both numerator and denominator by some number so as to make both the denominators equal:
\begin{equation*}
\frac{8}{3} - \frac{17}{12} = \frac{4 \times 8}{4 \times 3} - \frac{17}{12} = \frac{32}{12} - \frac{17}{12} = \frac{15}{12} = \frac{5}{4}
\end{equation*}
In this example it was sufficient to multiply the numerator and denominator of only one of the fractions by a number. In some other case it may be necessary to multiply both fractions:
\begin{equation*}
\frac{8}{3} - \frac{17}{4} = \frac{4 \times 8}{4 \times 3} - \frac{3 \times 17}{3 \times 4} = \frac{32}{12} - \frac{51}{12} = -\frac{19}{12}
\end{equation*}

\noindent \textbf{Point of notation:} The sum (difference) of fractions with the same denominator is often written as
$$
\frac{a}{c} \pm \frac{b}{c} = \frac{a \pm b}{c} 
$$

\subsubsection{Multiplication and division of fractions}

\noindent The multiplication of fractions comes down to multiplying the numerators and multiplying the denominators separately:
\begin{equation*}
\frac{a}{b} \times \frac{c}{d} = \frac{a \times c}{b \times d}
\end{equation*}
An example:
\begin{equation*}
\frac{4}{5} \times \frac{9}{2} = \frac{36}{10}
\end{equation*}
This fraction can still be simplified by dividing both numerator and denominator by $2$:
$$
\frac{36}{10} = \frac{18}{5}
$$
Dividing by a fraction is done by multiplying the first fraction by the inverse of the second fraction:
$$
\frac{a}{b} \div \frac{c}{d} = \frac{a}{b} \times \frac{d}{c} = \frac{a \times d}{b \times c}
$$
So that
\begin{equation*}
\frac{4}{5} \div \frac{9}{2} = \frac{4}{5} \times \frac{2}{9} = \frac{8}{45}
\end{equation*}
Dividing two numbers $a$ and $b$ is identical to multiplying $a$ with the inverse of $b$:
$$
\frac{a}{b} = a \times \frac{1}{b}
$$
This is consistent with the way multiplication of fractions was defined:
\begin{equation*}
5 \times \frac{1}{23} = \frac{5}{1} \times \frac{1}{23} = \frac{5 \times 1}{1 \times 23} = \frac{5}{23}
\end{equation*}
That is all there is to fractions...

\subsection{$\Sigma$: a shorthand notation for a sum}

\noindent The Greek capital letter $\Sigma$ (sigma) is used to shorten the writing of sums. Writing the sum of the 8 first (strictly) positive integers looks like
$$
1 + 2 + 3+ 4 + 5 + 6 + 7 + 8  = 36
$$
You will probably agree that this is a non-attractive, long-looking expression. The way a mathematician would write this is
$$
\sum\limits_{i = 1}^{8} i = 36
$$
This should be read as "the sum of $i$, for $i$ going from 1 to 8". In this expression the index $i$ is called a dummy index. It has no meaning as such. The information under the $\Sigma$ tells us that the first value that is plugged in to $i$, which in this example is 1. Then we write a plus and plug in the next value of $i$ which is 2. Then we write a plus...\\
And this is done until the last value of $i$ which is the value at the top of the $\Sigma$, in this case it's 8.\\

\noindent A couple of examples:
\begin{eqnarray*}
\sum\limits_{i = 0}^{5} \left( i\times i \right) &  = & 0\times 0 + 1 \times 1 + 2\times 2 + 3 \times 3 + 4 \times 4 + 5\times 5 = 55 \\
\sum\limits_{i = 3}^{6} \left(2i + 1\right) & = & 7 + 9 + 11 + 13 = 40 \\
\sum\limits_{i = 1}^{10} 3 & = & \underbrace{3 + 3 + 3 +\cdots + 3}_{10 \; terms} = 30\\ 
\sum\limits_{k = 2}^{6} \left( -\frac{2}{k}\right) &= & -\frac{2}{2} -\frac{2}{3} -\frac{2}{4}-\frac{2}{5} -\frac{2}{6} = -\frac{60 + 40 + 30 + 24 + 20}{60} = -\frac{174}{60} =  - \frac{29}{10}
\end{eqnarray*}
Notice in the third example the lack of "`$i$"' in the expression on the left hand side (LHS). The value to be summed remains 3 regardless of the value of $i$. Therefore we have to add a 3 for each value that $i$ can take, which is 10 in this case. In the last example we changed the letter we use as a dummy index from $i$ to $k$. This is always allowed as the dummy index has no meaning of its own. It is just a placeholder for whatever number we are going to plug in the expression.

\section{exponents, roots and absolute value}

\subsubsection{Positive integer exponents}

\noindent In the same way that a multiplication is a shorthand notation for a sum, an exponential is a shorthand notation for a multiplication:
\begin{eqnarray*}
6 \times 6 \times 6 & =  6^3 & =  216 \\
4 \times 4 \times 4 \times 4 & =  4^4 & =  256
\end{eqnarray*}
More generally, for any real number $a$ and any (strictly) positive integer $b$:
$$
a^b = \underbrace{a \times a \times \cdots \times a}_{b \; factors}
$$
where $a$ is called the base number and $b$ the exponent. We say "`$a$ to the power $b$"' or simply "`$a$ to the $b$-th"'.\\

\noindent From the definition it follows that $1^b$ = 1, no matter what the value of b is. If the exponent is 0 the result is by definition 1: $a^0 = 1$. The only exception to this is if $a$ itself is equal to 0 too, in which case it is simply not defined. So $0^0$ should never be used cause it is simply "bla".\\

\noindent If $b$ and $c$ are 2 positive integers than
$$
a^b \times a^c = \underbrace{a \times a \times \cdots \times a}_{b \; factors} \times \underbrace{a \times a \times \cdots \times a}_{c \; factors} = \underbrace{a \times a \times \cdots \times a}_{b+c \; factors} = a^{b+c}
$$
Another property which follows from the definition is
$$
\left( a^b \right)^c = \underbrace{a^b \times a^b \times \cdots \times a^b}_{c \; factors} = \overbrace{a^{b + b + \cdots + b}}^{c \; terms} = a^{b \times c} 
$$

\subsubsection{Negative integer exponents}

\noindent The negative (integer) exponent is defined as:
$$
a^{-b} = \frac{1}{a^b}
$$
So that
\begin{eqnarray*}
2^{-3} = \frac{1}{2^3} = \frac{1}{8} = 0,125 \\
\left( \frac{3}{5} \right)^{-4} = \frac{1}{\left( \frac{3}{5} \right)^4} = \left( \frac{5}{3} \right)^4 = \frac{5^4}{3^4} = \frac{625}{81}
\end{eqnarray*}

\subsubsection{Square root}

\noindent The square root is the inverse function of taking the square. This means
$$
a^2 = b \Leftrightarrow \sqrt{b} = a
$$
Examples are
\begin{eqnarray*}
\sqrt{16} = \pm 4 & \text{because} & \left( \pm 4 \right)^2 = 16 \\
 \sqrt{121} = \pm 11 & \text{because} & \left( \pm 11 \right)^2 = 121
\end{eqnarray*}
Notice that a square root has 2 results, because multiplying two positive reals as well as 2 negative reals gives a positive result. It is impossible to find a real number whose square is negative. Therefore the square root of a negative number doesn't exist\footnote{We will see later that there are indeed number whose square is negative, but these numbers do not belong to $\mathbb{R}$ and we can forget about them for the time being}.

\subsubsection{Other roots}

\noindent We saw that the square root is the inverse function of the square, i.e. taking the second power. More generally, the inverse function of taking any power is called a root. It is denoted as
$$
a^n = b \Leftrightarrow \sqrt[n]{b} = a
$$
In words: "`if the n-th power of a is b if and only if the \textbf{n-th root} of b is a"'. For example:
\begin{eqnarray*}
\sqrt[3]{8} =  2 & \text{because} & \left(  2 \right)^3 = 8 \\
 \sqrt[4]{81} = \pm 3 & \text{because} & \left( \pm 3 \right)^4 = 81
\end{eqnarray*}
Note that for even powers there are 2 results possible, because taking an even power of a base number or of the negative of the base number gives the same result.\\

\noindent The exponent we considered in the previous sections were integers. We now define the meaning of rational exponents. We take numbers $a,b$ and $c$ with the conditions that $c \neq 0 $, $b \in \mathbb{Z}$ and $a^b$ is non-negative if $c$ is even:  
$$
a^{\frac{b}{c}} =  \sqrt[c]{a^b}
$$
So that
\begin{eqnarray*}
4^{2,5}=4^{\frac{5}{2}} = \sqrt{4^5} = \sqrt{1024} = 32 \\
\left( -1,3 \right)^{-3,8} = \left( - \frac{13}{10} \right)^{-\frac{19}{5}} = \left( - \frac{10}{13} \right)^{\frac{19}{5}} = \left( \sqrt[5]{-\frac{10}{13}} \right)^{19} = -2,71
\end{eqnarray*}
The last example thus means that $\left( -2,71 \right)^{-\frac{1}{3,8}}$ = -1,3

\subsubsection{Absolute value}

\noindent The absolute value of a real number a is defined as:
$$
|a| = \left\{ \begin{array}{ll}
                  a \ \ \ \ \ \  \text{if } a \geq 0 \\
                  -a \ \ \ \ \ \  \text{if } a < 0 
                \end{array}
              \right.
$$
So if $a$ is positive or zero, then $|a| = a$, and if $a$ is negative then $|a| = -a$. Basically the absolute value is simply the real number where the minus sign, if there is one, is omitted:
\begin{eqnarray*}
 |19| = 19 \\
|-3| = 3 \\
|0| = 0
\end{eqnarray*}

\section{Order of operations}

\noindent How should an expression like the below be interpreted?
$$
2 + 3 \times 4
$$
\noindent If we were to read it from left to right the result would be $20$. We would start by adding $2$ and $3$ to get $5$. And multiply this $5$ by $4$ to get $20$. But this is wrong.\\

\noindent The convention in mathematics is to give precedence to multiplication and division, and only then do subtraction and addition. In the calculation above we thus first multiply $3$ by $4$ which is $12$. And then add $2$ to find the correct result of $14$.\\

\noindent Generally the convention is that operations have to be executed in the following order:
\begin{enumerate}
\item brackets
\item exponents 
\item multiplication and division
\item addition and subtraction
\end{enumerate}
\noindent Remember that roots, square or otherwise, are exponents and should therefore be treated as such.
\noindent Here are a couple of examples:
\begin{itemize}
\item $1 + 4 : 2 + 3 = 1 + 2 + 3 = 6$ 
\item $(1+4) : (2 + 3) = 5 : 5 = 1 $
\item $(1+4):2 + 3 = 5:2 + 3 = 2.5 + 3 = 5.5$
\item $\sqrt{5+4} : 2-3 = \sqrt{9} : 2 - 3 = 3 : 2 - 3 =  1.5 - 3 = -1.5$
\item $(3 + 8) \times (7-4) =  11 \times 3 = 33$
\item $(2+4)^3 = 6^3 = 216$
\item $(3 \times (12:4)^3 - 14) \times 2 = (3 \times 3^3 - 14) \times 2 = (81 - 14) \times 2 = 67 \times 2 = 134$
\item $\frac{1+2}{3+4} + 5 = \frac{3}{7} + 5 = 0.42857 + 5 = 5.42857$
\end{itemize}
\noindent In the last example we first worked out the sums in the fraction. Writing a fraction as in the last example should be interpreted as having brackets in the numerator and denominator:
$$
\frac{1+2}{3+4} + 5 = \frac{(1+2)}{(3+4)} + 5
$$

\section{Go and multiply!}

\noindent Keeping in mind the order of the operations, we do the following calculation:
\begin{eqnarray*}
(3+4)^2 & = & (3+4) \times (3+4) \\
& = & 7 \times 7 \\
& = & 49
\end{eqnarray*}
\noindent We get the same result if the brackets are worked out term by term:
\begin{equation*}
     \begin{tikzpicture}[>=stealth,baseline,anchor=base,inner sep=0pt]
       \matrix (foil) [matrix of math nodes,nodes={minimum height=0.5em}] {
         ( & 3 & + & 4 & ) & \times & ( & 3 & + & 4 & ) = 9 + 12 + 12 + 16 = 49\\
       };
       \path[->] ($(foil-1-2.north)+(0,1ex)$)   edge[red,bend left=45]    ($(foil-1-8.north)+(0,1ex)$)
                 ($(foil-1-2.north)+(0,1ex)$)   edge[green!40!black,bend left=60]  ($(foil-1-10.north)+(0,0.5ex)$)
                 ($(foil-1-4.north)+(0,0.5ex)$) edge[blue,bend left]      ($(foil-1-8.north)+(0,1ex)$)
                 ($(foil-1-4.north)+(0,0.5ex)$) edge[darkgray,bend left=45] ($(foil-1-10.north)+(0,0.5ex)$);
     \end{tikzpicture}
   \end{equation*}
\noindent This property is known as \textbf{distributivity of multiplication over addition} and holds generally:
\begin{equation*}
(a+b)^2=(a+b)  (a+b) = a^2 + a b + b  a + b^2 = a^2 + 2 a  b + b^2
\end{equation*}
Where we used the commutativity of the multiplication of numbers: $a\times b = b \times a$.\\

\noindent Use distributivity to convince yourself of the following identities:
\begin{center}
\begin{tabular}{ |c| }
  \hline
  $(a+b)^2 =  a^2 + 2ab + b^2$\\
  $(a-b)^2 =  a^2 - 2ab + b^2$\\
	$(a+b)(a-b) =   a^2 - b^2$\\
	$(a+b)^3 =  a^3 + 3 a^2b + 3 ab^2 + b^3$\\
	\hline
\end{tabular}
\end{center}